\section{Where the workbench meets the author's other workbenches}\label{sec:overlaps}

The owner notes that the Erdős--Straus and zeta sessions often ran at the same time, and that material crossed between them. The note \note{47\_} records every crossing that the audit found as a typed edge. Its edge types are:
\begin{itemize}
\item \emph{P}: a proof-level import;
\item \emph{D}: a dependency or citation;
\item \emph{I}: the same identity or object on both sides;
\item \emph{A}: an analogy only.
\end{itemize}
The table lists the edges that touch this workbench, and adds three that this reader found in the archive (marked ``new'').

\begin{center}\small
\begin{longtable}{@{}>{\raggedright\arraybackslash}p{0.13\linewidth}>{\raggedright\arraybackslash}p{0.44\linewidth}>{\raggedright\arraybackslash}p{0.07\linewidth}>{\raggedright\arraybackslash}p{0.28\linewidth}@{}}
\toprule
Edge & What crosses & Type & Status\\
\midrule
\endhead
NS $\to$ ES & The auxiliary torus endomorphism $J=\begin{pmatrix}3&1\\1&5\end{pmatrix}$ (determinant $14$) of the Navier--Stokes manuscript. It enters 21 statements of the supplement (§8) as an operator only. The core identity is $\{(u,v)\in G^2:u^3v=1=uv^5\}=\{(u,u^{-3}):u^{14}=1\}$ & P & audited (\note{21\_} §1); no fluid theorem is used\\
S$^6\to$ ES & The integral monodromies $A_1,A_2,A_\infty$ of the $(3,4,\infty)$ period system, in ES-09 & P & audited for odd levels (§\ref{ssec:es-s6})\\
ES $\leftrightarrow$ S$^6$ & The Niemeier lattice with root system $A_5^4D_4$ and glue index $72$ (ES result R10) & I & ES's glue code verified; the two codes not compared\\
Jacobian map & Alpöge's map $F:\C^3\to\C^3$ with $\det DF=-2$ and the three-point fibre over $(-\tfrac14,0,0)$ (archive §§4, 7). The same map is in the Yang--Mills workbench and the zeta reader. ES's four-dimensional extension $P$ ($\det DP=-2$) crosses to the zeta programme's conductor & P & $\det DF$, $\det DP$ and the fibre verified (zeta reader \cite{ZetaReader}, Lemma~6.6; \note{21\_}); the conductor maps located\\
ES $\leftrightarrow$ zeta & The split-zero semiring $G(R)=R\sqcup\{\tau\}$ with its Boolean support character (archive §21; ES R01) & I & audited (\note{21\_} §6)\\
ES $\leftrightarrow$ zeta & A finite Weyl phase-space basis & I & ES side verified on $C_2\times C_3$\\
zeta $\to$ ES & NG20--NG21, used by item 20; the SplitZero packages of 13--16 September & D & located\\
ES $\to$ zeta & The 488-page reader's source, vendored into the zeta repository & D & located\\
ES $\to$ zeta & The 124-page ES/Fable reader: solutions as binary quartics, then the Fable cover and a weighted conductor (§\ref{ssec:crosswalk}) & P & the quartic relation proved here (Proposition~\ref{prop:quartic}); the rest located\\
ES $\leftrightarrow$ Collatz & $3(4n+1)+1=4(3n+1)$, i.e.\ $\theta T=4\theta$ (ES R05, diagrams by u/CivQ17) & I & audited\\
Collatz $\leftrightarrow$ ES $\leftrightarrow$ zeta & The homogenisation $(x,h)\mapsto(ax+bh,h)$ (ES X4) & I & ES's diagram read\\
Busy Beaver (new) & Busy-Beaver extrema quoted by the archive ($S(4)=107$ and $\mathrm{Space}(5)=12289$; for the five-state values it cites the formal enumeration of the bbchallenge collaboration) select the shells of archive §§3 and 20; §25 maps extrema to shell carriers. The zeta reader's detectability argument uses a different machine for a different purpose & A & located; no mathematical dependence between the two\\
Connes (new) & Supplement §6 reads Connes' primitive intersections; archive §24 opens with Tomita--Takesaki transport. The zeta programme uses Connes' summation map and Meyer's quotient & A & located; shared names, no common statement found\\
Fable (new) & ``Fable'' names three different objects across the repositories: the three-dimensional Keller map, the $S^6$ construction in older files, and a conductor cover (\note{47a\_} §10). Archive §5 withdraws one Fable construction & --- & disambiguate at each use\\
ES $\leftrightarrow$ YM & The word ``gap'' (ES X7), and a proposed success projector & A & no mathematical content\\
\bottomrule
\end{longtable}
\end{center}

So the mathematical crossings are few and specific: the torus operator, the monodromy matrices, the Keller maps, the split-zero semiring, and two identities. Each is either audited or recorded as located. The remaining connections are shared vocabulary or proposals. The archive says the same of its late Fable, operator, Connes and moonshine material (p.~4): it is admitted only as a candidate interface, pending an explicit typed map and compatibility check.
